\chapter{Evaluación Económica Financiera}

\begingroup
\sloppy
\setlength{\emergencystretch}{3em}

\definecolor{EvalBlue}{HTML}{008CC7}
\definecolor{EvalBlueDark}{HTML}{006E9E}
\definecolor{EvalLight}{HTML}{CDEFF8}
\definecolor{EvalLighter}{HTML}{EAF8FC}
\definecolor{EvalTotal}{HTML}{00A6CF}
\definecolor{EvalGrid}{HTML}{1B4F63}

\newcommand{\mqiWhite}[1]{\textcolor{white}{\textbf{#1}}}
\newcommand{\mqiTableSetup}{\arrayrulecolor{EvalGrid}\renewcommand{\arraystretch}{1.12}\setlength{\tabcolsep}{3.5pt}}
\newcommand{\mqiTitle}[2]{\rowcolor{EvalBlueDark}\multicolumn{#1}{|c|}{\mqiWhite{#2}}\\ \hline}
\newcommand{\mqiHead}{\rowcolor{EvalBlue}}
\newcommand{\mqiA}{\rowcolor{EvalLight}}
\newcommand{\mqiB}{\rowcolor{EvalLighter}}
\newcommand{\mqiT}{\rowcolor{EvalTotal}}

\section{Criterio de evaluación y alcance económico}

El presente capítulo desarrolla la evaluación económica financiera de la solución tecnológica propuesta. El flujo fue construido siguiendo la lógica del libro de trabajo económico financiero utilizado como referencia docente y la explicación metodológica del documento de análisis económico financiero. Por ello, el análisis no se limita a una estimación simple de ingresos y egresos, sino que incorpora determinación de costos, planilla, gastos administrativos, gastos de venta, depreciaciones, amortizaciones, financiamiento, balance de apertura, presupuesto de caja, balance proyectado, punto de equilibrio, flujo de efectivo, VAN, TIR e indicadores de sensibilidad.

La unidad económica evaluada no es una venta aislada de hardware ni una aplicación móvil independiente. El proyecto se modela como un servicio tecnológico institucional para museos, compuesto por implementación, configuración técnica, estructuración curatorial, soporte recurrente, consumo de IA, analítica y mantenimiento. Esta definición recoge lo desarrollado en los capítulos previos: el mercado objetivo de museos peruanos, la cartera comercial, los paquetes de servicio, la arquitectura técnica, el uso de BLE/RAG y el diseño de despliegue.

\begin{table}[H]
	\centering
	\scriptsize
	\mqiTableSetup
	\resizebox{\textwidth}{!}{%
	\begin{tabular}{|c|p{5.1cm}|c|p{5.1cm}|}
		\hline
		\mqiTitle{4}{HOJAS DEL FLUJO DE CAJA DEL PROYECTO}
		\mqiHead \mqiWhite{N.} & \mqiWhite{Hoja} & \mqiWhite{N.} & \mqiWhite{Hoja} \\ \hline
		\mqiA 1 & INGRESOS & 14 & COMENTARIOS DE LAS RAZONES \\ \hline
		\mqiB 2 & COSTOS & 15 & PUNTO DE EQUILIBRIO \\ \hline
		\mqiA 3 & COSTO DE PRODUCCION UNITARIO & 16 & FLUJO DE EFECTIVO \\ \hline
		\mqiB 4 & PLANILLA & 17 & COSTO CAPITAL VAN TIR IR \\ \hline
		\mqiA 5 & GASTOS\_ADM\_VTAS\_FINANZAS & 18 & Sensibilidad Ingresos Bajos \\ \hline
		\mqiB 6 & DEPRECIACIONES AMORTIZACION & 19 & Sensibilidad Incremento Gastos \\ \hline
		\mqiA 7 & PLAN DE INVERSION & 20 & Sensibilidad Incremento Interes \\ \hline
		\mqiB 8 & BALANCE DE APERTURA & 21 & Riesgo no recuperacion CxC \\ \hline
		\mqiA 9 & COSTO DE VENTA & 22 & Riesgo Administrativo \\ \hline
		\mqiB 10 & ESTADO DE RESULTADOS & 23 & Flujo Optimista \\ \hline
		\mqiA 11 & PRESUPUESTO CAJA & 24 & Analisis Precio \\ \hline
		\mqiB 12 & BALANCE GENERAL PROYECTADO & 25 & Presupuesto Caja Bajando Precio \\ \hline
		\mqiA 13 & RAZONES FINANCIERAS & 26 & Flujo Efectivo Bajando Precio \\ \hline
	\end{tabular}%
	}
	\caption[Estructura del flujo financiero del proyecto]{Estructura del flujo financiero del proyecto. Fuente: elaboración propia con base en el archivo económico de referencia.}
	\label{tab:estructura-flujo-proyecto}
\end{table}

La evaluación se realiza en un horizonte de diez años, en soles nominales y sin IGV. El escenario base considera una captación progresiva de treinta y cinco museos al final del horizonte, con una mezcla comercial de 5 paquetes Básicos, 20 paquetes Estándar y 10 paquetes Avanzados. A diferencia de una evaluación idealizada, el modelo considera ventas al contado y al crédito, recuperación parcial de cuentas por cobrar dentro del año, pago diferido a proveedores, impuesto a la renta, deuda bancaria y variaciones por sensibilidad.

\begin{table}[H]
	\centering
	\footnotesize
	\mqiTableSetup
	\begin{tabular}{|p{8.2cm}|p{5.2cm}|}
		\hline
		\mqiTitle{2}{SUPUESTOS CENTRALES DEL FLUJO}
		\mqiHead \mqiWhite{Supuesto} & \mqiWhite{Valor adoptado} \\ \hline
		\mqiA Horizonte de evaluación & 10 años \\ \hline
		\mqiB Moneda & Soles nominales, sin IGV \\ \hline
		\mqiA Museos captados en escenario base & 35 museos acumulados al año 10 \\ \hline
		\mqiB Mezcla acumulada de paquetes & 5 Básicos, 20 Estándar y 10 Avanzados \\ \hline
		\mqiA Ajuste comercial sobre precios del Capítulo 5 & 20\% por formalización, garantía, gestión contractual y contingencia operativa \\ \hline
		\mqiB Servicios adicionales de despliegue & 8\% de la implementación anual \\ \hline
		\mqiA Política de cobranza & 40\% contado y 60\% crédito \\ \hline
		\mqiB Recuperación del crédito en el año & 80\% del crédito anual; el saldo queda como cuenta por cobrar \\ \hline
		\mqiA Política de compras directas & 30\% contado y 70\% crédito \\ \hline
		\mqiB Pago a proveedores en el año & 85\% del crédito anual; el saldo queda como cuenta por pagar \\ \hline
		\mqiA Impuesto a la renta & 29.5\% sobre utilidad antes de impuesto \\ \hline
		\mqiB Financiamiento & 40\% fondos propios y 60\% deuda \\ \hline
		\mqiA Tasa de deuda & 18\% anual, con dos años de gracia de capital \\ \hline
	\end{tabular}
	\caption[Supuestos centrales del flujo del proyecto]{Supuestos centrales del flujo económico financiero del proyecto. Fuente: elaboración propia.}
	\label{tab:supuestos-flujo-proyecto}
\end{table}

\section{Proyección de servicios e ingresos}

La capacidad de producción del servicio se expresa como número de museos implementados y número de museos activos con servicio recurrente. A diferencia del proyecto de referencia, donde la unidad productiva es un procedimiento médico, aquí la unidad de producción es un proyecto institucional de implementación y posterior operación. Cada museo requiere diagnóstico, instalación, configuración del backend, corpus curatorial, pruebas, capacitación y soporte.

El cronograma de captación comercial prioriza los paquetes Avanzados en los primeros años porque son los que permiten validar la propuesta en museos con mayor complejidad curatorial y mayor visibilidad institucional. Luego se incrementa la venta de paquetes Estándar y, desde el sexto año, se incorporan paquetes Básicos para instituciones de menor escala.

\begin{table}[H]
	\centering
	\scriptsize
	\mqiTableSetup
	\resizebox{\textwidth}{!}{%
	\begin{tabular}{|l|rrrrrrrrrr|r|}
		\hline
		\mqiTitle{12}{CRONOGRAMA DE CAPTACIÓN COMERCIAL}
		\mqiHead \mqiWhite{Paquete} & \mqiWhite{A1} & \mqiWhite{A2} & \mqiWhite{A3} & \mqiWhite{A4} & \mqiWhite{A5} & \mqiWhite{A6} & \mqiWhite{A7} & \mqiWhite{A8} & \mqiWhite{A9} & \mqiWhite{A10} & \mqiWhite{Total} \\ \hline
		\mqiA Básico & 0 & 0 & 0 & 0 & 0 & 1 & 1 & 1 & 1 & 1 & 5 \\ \hline
		\mqiB Estándar & 0 & 0 & 2 & 2 & 3 & 1 & 2 & 3 & 3 & 4 & 20 \\ \hline
		\mqiA Avanzado & 2 & 2 & 1 & 1 & 0 & 2 & 1 & 0 & 1 & 0 & 10 \\ \hline
		\mqiT \mqiWhite{Total anual} & \mqiWhite{2} & \mqiWhite{2} & \mqiWhite{3} & \mqiWhite{3} & \mqiWhite{3} & \mqiWhite{4} & \mqiWhite{4} & \mqiWhite{4} & \mqiWhite{5} & \mqiWhite{5} & \mqiWhite{35} \\ \hline
	\end{tabular}%
	}
	\caption[Cronograma de captación del proyecto]{Cronograma de captación comercial del escenario base. Fuente: hoja \texttt{INGRESOS}.}
	\label{tab:cronograma-captacion-proyecto}
\end{table}

Los precios económicos toman como base los paquetes definidos en el Capítulo 5 y agregan un factor comercial de 20\%. Este ajuste cubre formalización, garantía, gestión contractual, documentación, coordinación institucional y contingencias de operación. No se trata de inflar el precio arbitrariamente, sino de trasladar los precios técnicos a una operación empresarial completa.

\begin{table}[H]
	\centering
	\footnotesize
	\mqiTableSetup
	\resizebox{\textwidth}{!}{%
	\begin{tabular}{|l|r|r|r|r|p{3.0cm}|}
		\hline
		\mqiTitle{6}{VALORES ECONÓMICOS POR PAQUETE}
		\mqiHead \mqiWhite{Paquete} & \mqiWhite{Impl. Cap. 5} & \mqiWhite{Impl. flujo} & \mqiWhite{Recurrente Cap. 5} & \mqiWhite{Recurrente flujo} & \mqiWhite{Cobertura} \\ \hline
		\mqiA Básico & S/48\,000 & S/57\,600 & S/14\,000 & S/16\,800 & 3 a 5 salas \\ \hline
		\mqiB Estándar & S/96\,000 & S/115\,200 & S/29\,000 & S/34\,800 & 6 a 10 salas \\ \hline
		\mqiA Avanzado & S/185\,000 & S/222\,000 & S/52\,000 & S/62\,400 & 12 a 18 salas \\ \hline
	\end{tabular}%
	}
	\caption[Valores económicos por paquete del proyecto]{Valores económicos por paquete usados en el flujo de caja. Fuente: elaboración propia con base en el Capítulo 5.}
	\label{tab:precios-flujo-proyecto}
\end{table}

La Tabla~\ref{tab:ingresos-caja-proyecto} muestra la diferencia entre ventas devengadas e ingresos de efectivo. Esta separación es necesaria porque el servicio se vende a instituciones culturales, donde los pagos suelen depender de contratos, conformidad de entregables y tiempos administrativos. Por ello, el flujo no asume que toda venta se cobra en el mismo periodo.

\begin{table}[H]
	\centering
	\scriptsize
	\mqiTableSetup
	\resizebox{\textwidth}{!}{%
	\begin{tabular}{|l|rrrrrrrrrr|}
		\hline
		\mqiTitle{11}{VENTAS E INGRESOS DE EFECTIVO}
		\mqiHead \mqiWhite{Concepto} & \mqiWhite{A1} & \mqiWhite{A2} & \mqiWhite{A3} & \mqiWhite{A4} & \mqiWhite{A5} & \mqiWhite{A6} & \mqiWhite{A7} & \mqiWhite{A8} & \mqiWhite{A9} & \mqiWhite{A10} \\ \hline
		\mqiA Ventas totales & 541\,920 & 686\,722 & 853\,167 & 1\,023\,002 & 1\,056\,907 & 1\,598\,646 & 1\,719\,842 & 1\,813\,402 & 2\,399\,807 & 2\,570\,010 \\ \hline
		\mqiB Ingresos de efectivo & 476\,890 & 669\,345 & 833\,194 & 1\,002\,622 & 1\,052\,838 & 1\,533\,637 & 1\,705\,299 & 1\,802\,174 & 2\,329\,438 & 2\,549\,586 \\ \hline
		\mqiA Cuentas por cobrar al cierre & 65\,030 & 82\,407 & 102\,380 & 122\,760 & 126\,829 & 191\,837 & 206\,381 & 217\,608 & 287\,977 & 308\,401 \\ \hline
	\end{tabular}%
	}
	\caption[Ventas e ingresos de efectivo del proyecto]{Ventas, ingresos de efectivo y cuentas por cobrar del escenario base. Fuente: hoja \texttt{INGRESOS}.}
	\label{tab:ingresos-caja-proyecto}
\end{table}

\section{Determinación de costos}

Los costos del proyecto se organizan como en el documento de referencia: insumos y servicios directos, mano de obra de producción, costos indirectos de producción, gastos administrativos, gastos de venta y gastos financieros. En este caso, los insumos directos incluyen beacons BLE, nodos de borde, señalética, cableado, instalación, configuración, estructuración de corpus, pruebas, capacitación, documentación, nube, voz, RAG y soporte técnico.

El costo de venta se reconoce en el mismo año porque la propuesta presta un servicio por proyecto y no mantiene productos terminados. Por ello, el inventario final es cero y el costo de producción se transfiere íntegramente al costo de venta del periodo.

\begin{table}[H]
	\centering
	\scriptsize
	\mqiTableSetup
	\resizebox{\textwidth}{!}{%
	\begin{tabular}{|l|rrrrrrrrrr|}
		\hline
		\mqiTitle{11}{COSTO DE PRODUCCIÓN TRANSFERIDO}
		\mqiHead \mqiWhite{Concepto} & \mqiWhite{A1} & \mqiWhite{A2} & \mqiWhite{A3} & \mqiWhite{A4} & \mqiWhite{A5} & \mqiWhite{A6} & \mqiWhite{A7} & \mqiWhite{A8} & \mqiWhite{A9} & \mqiWhite{A10} \\ \hline
		\mqiA Insumos y servicios directos & 256\,760 & 299\,483 & 351\,260 & 401\,136 & 385\,007 & 613\,230 & 622\,589 & 622\,265 & 859\,276 & 879\,308 \\ \hline
		\mqiB Planilla de producción & 94\,043 & 143\,527 & 165\,929 & 174\,226 & 182\,937 & 207\,310 & 217\,675 & 273\,341 & 287\,008 & 301\,359 \\ \hline
		\mqiA Costos indirectos de servicio & 42\,675 & 52\,807 & 54\,344 & 55\,943 & 50\,478 & 46\,029 & 47\,828 & 49\,698 & 51\,643 & 53\,666 \\ \hline
		\mqiT \mqiWhite{Costo de producción} & \mqiWhite{393\,477} & \mqiWhite{495\,816} & \mqiWhite{571\,533} & \mqiWhite{631\,304} & \mqiWhite{618\,422} & \mqiWhite{866\,569} & \mqiWhite{888\,091} & \mqiWhite{945\,304} & \mqiWhite{1\,197\,928} & \mqiWhite{1\,234\,333} \\ \hline
	\end{tabular}%
	}
	\caption[Costo de producción del proyecto]{Costo de producción transferido al costo de venta. Fuente: hojas \texttt{COSTOS} y \texttt{COSTO DE VENTA}.}
	\label{tab:costos-produccion-proyecto}
\end{table}

La planilla se calculó con la estructura laboral empleada en el archivo de referencia: sueldo mensual, sueldo anual, gratificaciones, EsSalud, CTS, vacaciones y cesantía. La estructura inicia con un equipo reducido y se amplía de acuerdo con la cartera de museos, separando administración, producción y ventas.

\begin{table}[H]
	\centering
	\scriptsize
	\mqiTableSetup
	\resizebox{\textwidth}{!}{%
	\begin{tabular}{|l|rrrrrrrrrr|}
		\hline
		\mqiTitle{11}{PLANILLA POR ÁREA}
		\mqiHead \mqiWhite{Área} & \mqiWhite{A1} & \mqiWhite{A2} & \mqiWhite{A3} & \mqiWhite{A4} & \mqiWhite{A5} & \mqiWhite{A6} & \mqiWhite{A7} & \mqiWhite{A8} & \mqiWhite{A9} & \mqiWhite{A10} \\ \hline
		\mqiA Administración & 81\,504 & 85\,579 & 89\,858 & 94\,351 & 114\,294 & 120\,009 & 126\,009 & 132\,310 & 138\,925 & 145\,871 \\ \hline
		\mqiB Producción del servicio & 94\,043 & 143\,527 & 165\,929 & 174\,226 & 182\,937 & 207\,310 & 217\,675 & 273\,341 & 287\,008 & 301\,359 \\ \hline
		\mqiA Ventas & 15\,226 & 15\,987 & 16\,787 & 32\,852 & 34\,494 & 36\,219 & 38\,030 & 55\,158 & 57\,915 & 60\,811 \\ \hline
		\mqiT \mqiWhite{Total planilla} & \mqiWhite{190\,772} & \mqiWhite{245\,093} & \mqiWhite{272\,573} & \mqiWhite{301\,428} & \mqiWhite{331\,725} & \mqiWhite{363\,538} & \mqiWhite{381\,714} & \mqiWhite{460\,808} & \mqiWhite{483\,849} & \mqiWhite{508\,041} \\ \hline
	\end{tabular}%
	}
	\caption[Planilla por área del proyecto]{Planilla por área usada en el flujo de caja del proyecto. Fuente: hoja \texttt{PLANILLA}.}
	\label{tab:planilla-area-proyecto}
\end{table}

Los gastos administrativos y de ventas incluyen personal, servicios contables y legales, oficina, comunicaciones, prospección, viajes, marketing y comisiones comerciales. Los gastos financieros provienen del préstamo inicial de S/180\,000, con una tasa efectiva anual de 18\%, dos años de gracia de capital y amortización posterior hasta el año 8.

\section{Inversión, depreciación y financiamiento}

La inversión total requerida asciende a S/300\,000. Esta inversión no corresponde a maquinaria pesada, sino a la estructura mínima para iniciar, vender, implementar y sostener el servicio tecnológico propuesto. Los activos fijos comprenden laptops, equipos de prueba, servidor/NAS, red local, herramientas de instalación y mobiliario. Los activos nominales incluyen constitución, formalización, desarrollo base de app/API/RAG, documentación, manuales, marca, preventa y dossier comercial. El capital de trabajo cubre el ciclo inicial de diagnóstico, contratación de insumos, entrega, conformidad y cobro institucional.

\begin{table}[H]
	\centering
	\footnotesize
	\mqiTableSetup
	\resizebox{\textwidth}{!}{%
	\begin{tabular}{|p{6.2cm}|r|r|r|}
		\hline
		\mqiTitle{4}{PLAN DE INVERSIÓN Y FINANCIAMIENTO}
		\mqiHead \mqiWhite{Concepto} & \mqiWhite{Monto} & \mqiWhite{Fondos propios} & \mqiWhite{Financiamiento} \\ \hline
		\mqiA Activos fijos tangibles & 118\,000 & 23\,600 & 94\,400 \\ \hline
		\mqiB Capital de trabajo inicial & 50\,000 & 50\,000 & 0 \\ \hline
		\mqiA Activos nominales e intangibles & 132\,000 & 46\,400 & 85\,600 \\ \hline
		\mqiT \mqiWhite{Total inversión} & \mqiWhite{300\,000} & \mqiWhite{120\,000} & \mqiWhite{180\,000} \\ \hline
		\mqiB Participación & 100\% & 40\% & 60\% \\ \hline
	\end{tabular}%
	}
	\caption[Plan de inversión del proyecto]{Plan de inversión y fuentes de financiamiento del proyecto. Fuente: hoja \texttt{PLAN DE INVERSION}.}
	\label{tab:plan-inversion-proyecto}
\end{table}

\begin{table}[H]
	\centering
	\footnotesize
	\mqiTableSetup
	\begin{tabular}{|p{8.5cm}|r|}
		\hline
		\mqiTitle{2}{DETALLE PRINCIPAL DE ACTIVOS}
		\mqiHead \mqiWhite{Concepto} & \mqiWhite{Monto} \\ \hline
		\mqiA Laptops y equipos de prueba & 48\,000 \\ \hline
		\mqiB Servidor/NAS y red local & 32\,000 \\ \hline
		\mqiA Herramientas de instalación y medición & 20\,000 \\ \hline
		\mqiB Mobiliario y oficina & 18\,000 \\ \hline
		\mqiT \mqiWhite{Total activos fijos} & \mqiWhite{118\,000} \\ \hline
		\mqiA Constitución y formalización & 8\,000 \\ \hline
		\mqiB Desarrollo base app/API/RAG & 90\,000 \\ \hline
		\mqiA Documentación y manuales & 22\,000 \\ \hline
		\mqiB Marca, preventa y dossier comercial & 12\,000 \\ \hline
		\mqiT \mqiWhite{Total activos nominales} & \mqiWhite{132\,000} \\ \hline
	\end{tabular}
	\caption[Detalle de activos del proyecto]{Detalle principal de activos fijos y nominales del proyecto. Fuente: hoja \texttt{PLAN DE INVERSION}.}
	\label{tab:detalle-activos-proyecto}
\end{table}

El balance de apertura se construye directamente desde el plan de inversión. La caja inicial corresponde al capital de trabajo aprobado, los activos fijos y diferidos se registran por su valor de adquisición, la deuda de largo plazo corresponde al financiamiento bancario y el patrimonio corresponde al aporte de los socios.

\begin{table}[H]
	\centering
	\footnotesize
	\mqiTableSetup
	\begin{tabular}{|p{5.6cm}|r|p{5.6cm}|r|}
		\hline
		\mqiTitle{4}{BALANCE DE APERTURA}
		\mqiHead \mqiWhite{Activos} & \mqiWhite{Monto} & \mqiWhite{Pasivo y patrimonio} & \mqiWhite{Monto} \\ \hline
		\mqiA Caja y bancos & 50\,000 & Deuda de largo plazo & 180\,000 \\ \hline
		\mqiB Activo fijo & 118\,000 & Capital social aportado & 120\,000 \\ \hline
		\mqiA Activo diferido & 132\,000 & & \\ \hline
		\mqiT \mqiWhite{Total activos} & \mqiWhite{300\,000} & \mqiWhite{Total pasivo y patrimonio} & \mqiWhite{300\,000} \\ \hline
	\end{tabular}
	\caption[Balance de apertura del proyecto]{Balance de apertura del proyecto. Fuente: hoja \texttt{BALANCE DE APERTURA}.}
	\label{tab:balance-apertura-proyecto}
\end{table}

Las depreciaciones y amortizaciones se calculan en una hoja independiente, de forma análoga al archivo de referencia. La depreciación anual de activos tangibles asciende a S/23\,958 durante los primeros cuatro años, baja a S/12\,078 en el quinto año y luego a S/1\,782 anuales. La amortización de activos nominales es de S/26\,400 anuales durante los primeros cinco años. Estos importes se trasladan a producción y administración según la asignación contable del flujo.

\section{Estados financieros proyectados}

El estado de resultados muestra pérdidas netas en los dos primeros años, lo cual es consistente con una solución tecnológica institucional que todavía absorbe estructura, deuda y costos de implementación. A partir del tercer año la utilidad neta se vuelve positiva y se acelera cuando la base de museos activos genera ingresos recurrentes.

\begin{table}[H]
	\centering
	\scriptsize
	\mqiTableSetup
	\resizebox{\textwidth}{!}{%
	\begin{tabular}{|l|rrrrrrrrrr|}
		\hline
		\mqiTitle{11}{ESTADO DE RESULTADOS RESUMIDO}
		\mqiHead \mqiWhite{Concepto} & \mqiWhite{A1} & \mqiWhite{A2} & \mqiWhite{A3} & \mqiWhite{A4} & \mqiWhite{A5} & \mqiWhite{A6} & \mqiWhite{A7} & \mqiWhite{A8} & \mqiWhite{A9} & \mqiWhite{A10} \\ \hline
		\mqiA Ventas totales & 541\,920 & 686\,722 & 853\,167 & 1\,023\,002 & 1\,056\,907 & 1\,598\,646 & 1\,719\,842 & 1\,813\,402 & 2\,399\,807 & 2\,570\,010 \\ \hline
		\mqiB Costo de venta & 393\,477 & 495\,816 & 571\,533 & 631\,304 & 618\,422 & 866\,569 & 888\,091 & 945\,304 & 1\,197\,928 & 1\,234\,333 \\ \hline
		\mqiA Utilidad bruta & 148\,443 & 190\,905 & 281\,634 & 391\,698 & 438\,485 & 732\,077 & 831\,751 & 868\,097 & 1\,201\,879 & 1\,335\,677 \\ \hline
		\mqiB Gastos de operación & 178\,261 & 187\,997 & 198\,568 & 224\,756 & 243\,878 & 235\,840 & 248\,238 & 275\,625 & 301\,343 & 317\,192 \\ \hline
		\mqiA Utilidad de operación & -29\,818 & 2\,908 & 83\,067 & 166\,942 & 194\,607 & 496\,237 & 583\,512 & 592\,472 & 900\,536 & 1\,018\,486 \\ \hline
		\mqiB Gastos financieros & 32\,400 & 32\,400 & 32\,400 & 28\,969 & 24\,919 & 20\,141 & 14\,503 & 7\,850 & 0 & 0 \\ \hline
		\mqiA Impuesto a la renta & 0 & 0 & 14\,947 & 40\,702 & 50\,058 & 140\,448 & 167\,858 & 172\,463 & 265\,658 & 300\,453 \\ \hline
		\mqiT \mqiWhite{Utilidad neta} & \mqiWhite{-62\,218} & \mqiWhite{-29\,492} & \mqiWhite{35\,720} & \mqiWhite{97\,271} & \mqiWhite{119\,630} & \mqiWhite{335\,648} & \mqiWhite{401\,151} & \mqiWhite{412\,158} & \mqiWhite{634\,878} & \mqiWhite{718\,032} \\ \hline
	\end{tabular}%
	}
	\caption[Estado de resultados del proyecto]{Estado de resultados resumido del escenario base. Fuente: hoja \texttt{ESTADO DE RESULTADOS}.}
	\label{tab:estado-resultados-proyecto}
\end{table}

El presupuesto de caja permite observar la liquidez real del proyecto. Aunque el estado de resultados muestra pérdidas contables en los primeros años, el flujo mantiene saldo positivo gracias al capital de trabajo inicial y a la cobranza progresiva. El año más ajustado es el primero, con un saldo final de caja de S/69; por tanto, la operación no presenta déficit en el escenario base, pero sí exige disciplina en anticipos, cobranza y control de pagos.

\begin{table}[H]
	\centering
	\scriptsize
	\mqiTableSetup
	\resizebox{\textwidth}{!}{%
	\begin{tabular}{|l|rrrrrrrrrr|}
		\hline
		\mqiTitle{11}{PRESUPUESTO DE CAJA}
		\mqiHead \mqiWhite{Concepto} & \mqiWhite{A1} & \mqiWhite{A2} & \mqiWhite{A3} & \mqiWhite{A4} & \mqiWhite{A5} & \mqiWhite{A6} & \mqiWhite{A7} & \mqiWhite{A8} & \mqiWhite{A9} & \mqiWhite{A10} \\ \hline
		\mqiA Total ingresos & 476\,890 & 669\,345 & 833\,194 & 1\,002\,622 & 1\,052\,838 & 1\,533\,637 & 1\,705\,299 & 1\,802\,174 & 2\,329\,438 & 2\,549\,586 \\ \hline
		\mqiB Total egresos & 526\,820 & 661\,370 & 765\,770 & 866\,876 & 917\,681 & 1\,178\,185 & 1\,325\,477 & 1\,438\,503 & 1\,645\,066 & 1\,813\,297 \\ \hline
		\mqiA Flujo del periodo & -49\,931 & 7\,976 & 67\,424 & 135\,746 & 135\,157 & 355\,452 & 379\,821 & 363\,672 & 684\,373 & 736\,288 \\ \hline
		\mqiT \mqiWhite{Saldo final de caja} & \mqiWhite{69} & \mqiWhite{8\,045} & \mqiWhite{75\,469} & \mqiWhite{211\,215} & \mqiWhite{346\,372} & \mqiWhite{701\,825} & \mqiWhite{1\,081\,646} & \mqiWhite{1\,445\,318} & \mqiWhite{2\,129\,690} & \mqiWhite{2\,865\,979} \\ \hline
	\end{tabular}%
	}
	\caption[Presupuesto de caja del proyecto]{Presupuesto de caja del escenario base. Fuente: hoja \texttt{PRESUPUESTO CAJA}.}
	\label{tab:presupuesto-caja-proyecto}
\end{table}

El balance general proyectado cuadra en todos los años del modelo. Al año 10, los activos totales alcanzan S/3\,175\,560, de los cuales S/2\,865\,979 corresponden a caja y bancos y S/308\,401 a cuentas por cobrar. El pasivo total llega a S/392\,781, compuesto por cuentas por pagar e impuesto por pagar, mientras que el patrimonio alcanza S/2\,782\,779. Esta evolución confirma que el valor económico se concentra en la caja generada por la operación y en la acumulación de utilidades.

\section{Flujo de efectivo e indicadores financieros}

El flujo de efectivo económico considera la inversión inicial, el flujo anual de caja, la liquidación de cuentas por cobrar al cierre del horizonte, el valor de salvamento y la cancelación del pasivo circulante final. La inversión inicial se registra como salida en el año 0; luego se evalúan los flujos netos descontados con la tasa de corte calculada en el propio libro de trabajo.

\begin{table}[H]
	\centering
	\scriptsize
	\mqiTableSetup
	\resizebox{\textwidth}{!}{%
	\begin{tabular}{|l|r|rrrrrrrrrr|}
		\hline
		\mqiTitle{12}{FLUJO NETO DE EFECTIVO PARA EVALUACIÓN}
		\mqiHead \mqiWhite{Concepto} & \mqiWhite{A0} & \mqiWhite{A1} & \mqiWhite{A2} & \mqiWhite{A3} & \mqiWhite{A4} & \mqiWhite{A5} & \mqiWhite{A6} & \mqiWhite{A7} & \mqiWhite{A8} & \mqiWhite{A9} & \mqiWhite{A10} \\ \hline
		\mqiA Flujo para evaluación & -300\,000 & -49\,931 & 7\,976 & 67\,424 & 135\,746 & 135\,157 & 355\,452 & 379\,821 & 363\,672 & 684\,373 & 686\,909 \\ \hline
	\end{tabular}%
	}
	\caption[Flujo neto de efectivo del proyecto]{Flujo neto de efectivo usado para VAN, TIR e índice de rentabilidad. Fuente: hoja \texttt{FLUJO DE EFECTIVO}.}
	\label{tab:flujo-neto-proyecto}
\end{table}

La tasa de corte se calcula con la misma lógica del archivo de referencia. Primero se obtiene el costo promedio ponderado del capital: 40\% de fondos propios con costo de 9\% y 60\% de deuda con costo de 18\%. Esto produce un costo de capital de 14.40\%. Luego se incorpora una inflación promedio de 3.55\% y una prima de riesgo equivalente a 0.51\%. La tasa de corte resultante es 18.46\%.

\begin{table}[H]
	\centering
	\footnotesize
	\mqiTableSetup
	\begin{tabular}{|p{5.0cm}|r|r|r|}
		\hline
		\mqiTitle{4}{CÁLCULO DE LA TASA DE CORTE}
		\mqiHead \mqiWhite{Componente} & \mqiWhite{Monto / base} & \mqiWhite{Costo / tasa} & \mqiWhite{Aporte} \\ \hline
		\mqiA Fondos propios & S/120\,000 & 9.00\% & 3.60\% \\ \hline
		\mqiB Financiamiento & S/180\,000 & 18.00\% & 10.80\% \\ \hline
		\mqiT \mqiWhite{Costo de capital ponderado} & \mqiWhite{S/300\,000} & \mqiWhite{} & \mqiWhite{14.40\%} \\ \hline
		\mqiA Inflación promedio proyectada & -- & 3.55\% & 3.55\% \\ \hline
		\mqiB Riesgo de inversión & 14.40\% $\times$ 3.55\% & -- & 0.51\% \\ \hline
		\mqiT \mqiWhite{Tasa de corte} & \mqiWhite{} & \mqiWhite{} & \mqiWhite{18.46\%} \\ \hline
	\end{tabular}
	\caption[Cálculo de tasa de corte del proyecto]{Cálculo de la tasa de corte usada para VAN y TIR. Fuente: hoja \texttt{COSTO CAPITAL VAN TIR IR}.}
	\label{tab:tasa-corte-proyecto}
\end{table}

Con la tasa de corte de 18.46\%, el proyecto genera un VAN positivo de S/444\,585, una TIR de 33.99\% y un índice de rentabilidad de 2.48. El periodo de recuperación descontado es 6.35 años. En consecuencia, el proyecto se acepta en el escenario base, porque el VAN es mayor que cero, la TIR supera la tasa de corte y el índice de rentabilidad es mayor que uno.

\begin{table}[H]
	\centering
	\footnotesize
	\mqiTableSetup
	\begin{tabular}{|p{5.9cm}|r|p{5.7cm}|}
		\hline
		\mqiTitle{3}{INDICADORES FINANCIEROS}
		\mqiHead \mqiWhite{Indicador} & \mqiWhite{Resultado} & \mqiWhite{Lectura} \\ \hline
		\mqiA Tasa de corte & 18.46\% & Tasa mínima exigida para aceptar el proyecto \\ \hline
		\mqiB Valor actual neto (VAN) & S/444\,585 & Genera valor económico positivo \\ \hline
		\mqiA Tasa interna de retorno (TIR) & 33.99\% & Supera la tasa de corte en 15.53 puntos porcentuales \\ \hline
		\mqiB Índice de rentabilidad (IR) & 2.48 & Por cada sol invertido se obtiene un valor presente superior al desembolso inicial \\ \hline
		\mqiA Recuperación descontada & 6.35 años & La inversión se recupera durante el séptimo año de operación \\ \hline
		\mqiB Resultado & Proyecto se acepta & Cumple VAN positivo, TIR mayor que tasa de corte e IR mayor que uno \\ \hline
	\end{tabular}
	\caption[Indicadores financieros del proyecto]{Indicadores financieros del escenario base. Fuente: hoja \texttt{COSTO CAPITAL VAN TIR IR}.}
	\label{tab:indicadores-financieros-proyecto}
\end{table}

\section{Punto de equilibrio y lectura operativa}

El punto de equilibrio permite evaluar si las ventas anuales cubren los costos variables y fijos. En los dos primeros años el punto de equilibrio supera las ventas, lo que refleja el peso de la estructura inicial y de la deuda. A partir del tercer año, las ventas superan el punto de equilibrio y el margen mejora de manera sostenida. En el año 10, el punto de equilibrio equivale a 28.6\% de las ventas, lo cual muestra una operación más estable.

\begin{table}[H]
	\centering
	\scriptsize
	\mqiTableSetup
	\resizebox{\textwidth}{!}{%
	\begin{tabular}{|l|rrrrrrrrrr|}
		\hline
		\mqiTitle{11}{PUNTO DE EQUILIBRIO}
		\mqiHead \mqiWhite{Concepto} & \mqiWhite{A1} & \mqiWhite{A2} & \mqiWhite{A3} & \mqiWhite{A4} & \mqiWhite{A5} & \mqiWhite{A6} & \mqiWhite{A7} & \mqiWhite{A8} & \mqiWhite{A9} & \mqiWhite{A10} \\ \hline
		\mqiA Punto de equilibrio en soles & 694\,448 & 762\,750 & 733\,133 & 724\,811 & 709\,553 & 655\,754 & 645\,076 & 698\,276 & 722\,579 & 735\,109 \\ \hline
		\mqiB Porcentaje sobre ventas & 128.1\% & 111.1\% & 85.9\% & 70.9\% & 67.1\% & 41.0\% & 37.5\% & 38.5\% & 30.1\% & 28.6\% \\ \hline
	\end{tabular}%
	}
	\caption[Punto de equilibrio del proyecto]{Punto de equilibrio del escenario base. Fuente: hoja \texttt{PUNTO DE EQUILIBRIO}.}
	\label{tab:punto-equilibrio-proyecto}
\end{table}

La lectura operativa es clara: el proyecto necesita sostener la captación de paquetes Estándar y Avanzados para absorber costos fijos y mantener una base recurrente suficiente. La rentabilidad no proviene únicamente de vender implementaciones nuevas, sino de acumular museos activos que pagan soporte, IA, analítica y mantenimiento.

\section{Análisis de sensibilidad}

El análisis de sensibilidad permite identificar las condiciones bajo las cuales el proyecto deja de ser atractivo. El escenario base es viable, pero no debe interpretarse como garantía automática. La reducción de ingresos de 20\% vuelve negativo el VAN; la no recuperación del 20\% de cuentas por cobrar también vuelve inviable el proyecto; y una reducción de precios de 15\% destruye valor. En cambio, el incremento de costos y gastos de 10\%, el aumento de tasa de interés a 24\%, el incremento de gastos administrativos de 15\% y el escenario de precio -10\% aún mantienen viabilidad, aunque con menor holgura.

\begin{table}[H]
	\centering
	\footnotesize
	\mqiTableSetup
	\resizebox{\textwidth}{!}{%
	\begin{tabular}{|p{5.0cm}|r|r|r|p{4.7cm}|}
		\hline
		\mqiTitle{5}{ANÁLISIS DE SENSIBILIDAD}
		\mqiHead \mqiWhite{Escenario} & \mqiWhite{VAN} & \mqiWhite{TIR} & \mqiWhite{IR} & \mqiWhite{Resultado} \\ \hline
		\mqiA Base & 444\,585 & 33.99\% & 2.48 & Viable \\ \hline
		\mqiB Ingresos -20\% & -324\,974 & 6.29\% & -0.08 & No viable \\ \hline
		\mqiA Costos y gastos +10\% & 141\,276 & 23.26\% & 1.47 & Viable, con menor margen \\ \hline
		\mqiB Interés del financiamiento 24\% & 261\,962 & 32.64\% & 1.87 & Viable, pero con mayor tasa de corte \\ \hline
		\mqiA No recuperación CxC 20\% & -126\,994 & 13.41\% & 0.58 & No viable \\ \hline
		\mqiB Gastos administrativos +15\% & 378\,808 & 31.49\% & 2.26 & Viable \\ \hline
		\mqiA Precio -10\% & 69\,017 & 20.93\% & 1.23 & Viable, con liquidez ajustada \\ \hline
		\mqiB Precio -15\% & -123\,659 & 13.96\% & 0.59 & No viable \\ \hline
		\mqiA Optimista & 793\,866 & 45.62\% & 3.65 & Viable con amplia holgura \\ \hline
	\end{tabular}%
	}
	\caption[Análisis de sensibilidad del proyecto]{Análisis de sensibilidad financiero del proyecto. Fuente: hojas de sensibilidad del flujo de caja.}
	\label{tab:sensibilidad-proyecto}
\end{table}

La sensibilidad muestra tres riesgos críticos. El primero es comercial: si la captación o el precio efectivo caen demasiado, el proyecto no cubre el costo de capital. El segundo es financiero-operativo: la no recuperación de cuentas por cobrar afecta fuertemente el valor del proyecto, por lo que los contratos deben exigir anticipos, hitos de pago y conformidad documentada. El tercero es de alcance: una reducción de precio solo es sostenible si se reduce el alcance técnico, el consumo de IA o el soporte incluido. En particular, el precio -10\% aún resulta viable, pero el propio presupuesto de caja de ese escenario muestra una necesidad temporal de liquidez cercana a S/111\,946; por ello, no debería aplicarse sin capital de trabajo adicional o condiciones de cobro más estrictas.

\section{Conclusión financiera}

La evaluación económica financiera muestra que el proyecto es viable en el escenario base. La propuesta requiere una inversión inicial de S/300\,000, financiada con S/120\,000 de fondos propios y S/180\,000 de deuda. Con una cartera acumulada de 35 museos en diez años, ventas institucionales al contado y al crédito, costos de implementación, soporte recurrente, planilla, gastos, depreciaciones, amortizaciones, impuestos y deuda, el flujo obtiene un VAN de S/444\,585, una TIR de 33.99\%, un índice de rentabilidad de 2.48 y una recuperación descontada de 6.35 años.

La conclusión no es que la solución sea rentable bajo cualquier condición, sino que es financieramente defendible si se gestiona como un servicio institucional completo. La propuesta debe monetizar implementación, corpus curatorial, soporte, IA, analítica y mantenimiento, evitando competir como simple venta de beacons o aplicación genérica. Asimismo, la operación debe proteger la caja mediante anticipos, control de cuentas por cobrar y límites claros de uso de IA/voz.

En consecuencia, el proyecto se acepta bajo el escenario base. Sin embargo, su viabilidad queda condicionada a mantener precios consistentes con el alcance, priorizar contratos Estándar y Avanzados, controlar costos recurrentes de nube e IA, y asegurar mecanismos de cobranza contractual. Bajo esas condiciones, el prototipo técnico desarrollado en los capítulos previos puede evolucionar hacia una operación sostenible para museos peruanos.

\arrayrulecolor{black}
\endgroup
